%!TEX program = xelatex
%!TEX encoding = ISO-8859-1
\documentclass[12pt]{article}

\usepackage{amssymb}
\usepackage{amsmath,amsfonts,amsthm,amstext,mathabx}
\usepackage[brazil]{babel}
%\usepackage[latin1]{inputenc}
\usepackage{graphicx}
\graphicspath{{/home/jfreitas/Dropbox/imagens-latex/}{/ArquivosUbuntu/Dropbox/imagens-latex/}{D:/Dropbox/imagens-latex/}}
\usepackage{enumitem}
\usepackage{multicol}
\usepackage[all]{xy}

\setlength{\topmargin}{-1.0in}
\setlength{\oddsidemargin}{0in}
\setlength{\textheight}{10.1in}
\setlength{\textwidth}{6.5in}
\setlength{\baselineskip}{12mm}

\newcounter{exercicios}
\renewcommand{\theexercicios}{\arabic{exercicios}}

\newenvironment{questao}[1]{%
\refstepcounter{exercicios}%
\label{#1}%
\noindent\textbf{Exerc{\'\i}cio {\theexercicios}:}%
}%
{}%

% \newcounter{exercicios}
% \setcounter{exercicios}{0}
% \newcommand{\questao}{
% \addtocounter{exercicios}{1}
% \noindent{\bf Exerc{\'\i}cio \arabic{exercicios}: }}

\newcommand{\equi}{\Leftrightarrow}
\newcommand{\bic}{\leftrightarrow}
\newcommand{\cond}{\rightarrow}
\newcommand{\impl}{\Rightarrow}
\newcommand{\nao}{\sim}
\newcommand{\sub}{\subseteq}
\newcommand{\e}{\ \wedge\ }
\newcommand{\ou}{\ \vee\ }
\newcommand{\vaz}{\emptyset}
\newcommand{\nsub}{\nsubseteq}
\renewcommand{\sin}{{\rm sen\,}}

\newcommand{\n}{\mathbb{N}}
\newcommand{\z}{\mathbb{Z}}
\newcommand{\q}{\mathbb{Q}}
\newcommand{\im}{{\rm Im\,}}
\newcommand{\aut}{{\rm Aut\,}}
\newcommand{\cp}[1]{\mathbb{#1}}
\newcommand{\real}{\mathbb{R}}
\newcommand{\complex}{\mathbb{C}}
\newcommand{\lap}[1]{\mathcal{L}\left\{#1\right\}}
\newcommand{\lapi}[1]{\mathcal{L}^{-1}\left\{#1\right\}}
\newcommand{\se}[1]{\displaystyle\sum_{n = 1}^\infty{#1}}
\newcommand{\dlim}[2]{\displaystyle\lim_{#1\rightarrow #2}}
\newcommand{\slim}{\displaystyle\lim_{n \rightarrow \infty}}
\newcommand{\seq}[1]{\{{#1_n\}}}
\newcommand{\seg}[1]{\displaystyle\sum_{n = 1}^\infty{#1_n}}
\newcommand{\sei}[2]{\displaystyle\sum_{#1}^\infty{#2}}
\newcommand{\sepc}[3]{\displaystyle\sum_{#1}^\infty{#2(x - #3)^n}}
\newcommand{\imp}[3]{\displaystyle\int_{#1}^{+\infty}{#3}{d #2}}
\newcommand{\dint}[4]{\displaystyle\int_{#1}^{#2}{#4}{d#3}}
\newcommand{\inti}[2]{\displaystyle\int{#1}{d#2}}
\newcommand{\norma}[1]{\left\lVert#1\right\rVert}
\newcommand{\flim}[1]{\displaystyle\lim_{#1\rightarrow \infty}}
\renewcommand{\sin}{{\rm sen\,}}
\renewcommand{\tan}{{\rm tg\,}}
\renewcommand{\csc}{{\rm cossec\,}}
\renewcommand{\cot}{{\rm cotg\,}}
\renewcommand{\sinh}{{\rm senh\,}}
\newcommand\T{\rule{0pt}{2.6ex}} 
\newcommand{\vesp}{\vspace{0.2cm}}
\newcommand{\subne}{\subsetneqq}


\newcommand{\compcent}[1]{\vcenter{\hbox{$#1\circ$}}}
\newcommand{\comp}{\mathbin{\mathchoice
{\compcent\scriptstyle}{\compcent\scriptstyle}
{\compcent\scriptscriptstyle}{\compcent\scriptscriptstyle}}}

\begin{document}

\pagestyle{empty}

\begin{figure}[h]
        \begin{minipage}[c]{1.7cm}
        \includegraphics[width=1.7cm]{unb.pdf}
        \end{minipage}%
        \hspace{0pt}
        \begin{minipage}[c]{4in}
          {Universidade de Bras{\'\i}lia} \\
          {Departamento de Matem{\'a}tica}
\end{minipage}
\end{figure}
\vspace{-1cm}\hrule


\begin{center}
{\Large\bf {\'A}lgebra 1 - Turma D -- 2$^{o}$/2017} \\ \vspace{9pt} {\large\bf
  $3^{\underline{o}}$ Teste Módulo 3 - Resolu\c{c}\~ao}\\
\vspace{9pt} Prof. Jos{\'e} Ant{\^o}nio O. Freitas
\end{center}
\hrule

\vspace{.6cm}

\questao Em $\z \times\z$ considere a operação
\[
	(x, y) + (z, t) = (x + z, y + t)
\]
para todos $(x,y)$, $(z,t) \in \z\times\z$. Mostre que $(\z\times\z, +)$ é um grupo e que este grupo é abeliano.

\noindent\textbf{Solu\c{c}\~ao:}

Primeiro dados $(x,y)$, $(z,t) \in \z\times\z$ temos
\[
	(x,y) + (z,t) = (x+z,y+t) = (z+x,y+t) = (z,t) + (x,y).
\]
Assim a operação em $\z\times\z$ é comutativa.

Sejam $(x,y)$, $(z,t)$ e $(r,s) \in \z\times\z$. Então
\begin{align*}
	[(x,y) + (z,t)] + (r,s) &= (x + z, y+t) + (r,s) = (x+y+r,y+t+s)\\
	(x,y) + [(z,t)+(r,s)] &= (x,y) + (z+r,t+s) = (x+z+r,y+t+s).
\end{align*}
Logo $[(x,y)+(z,t)] + (r,s) = (x,y) + [(z,t) + (r,s)]$.

Para todo $(x,y) \in \z\times\z$ temos
\[
	(x,y) + (0,0) = (x+0,y+0) = (x,y).
\]
Assim $(0,0)$ é o elemento neutro da operação em $\z\times\z$.

Agora dado $(x,y) \in \z\times\z$ tome $(-x,-y) \in \z\times\z$. Então
\[
	(x,y) + (-x,-y) = (x+(-x),y+(-y)) = (0,0).
\]
Daí $(-x,-y)$ é o oposto de $(x,y)$ em $\z\times\z$.

Portanto $(\z\times\z,+)$ é um grupo abeliano.

\vspace{1cm}

\questao Seja $f : \z\times\z \to \z_7\times\z_7$ definida por
\[
	f(x,y) = (\overline{x-y},\overline{0}).
\]
Mostre que $f$ é um homomorfismo de grupos e calcule $\ker(f)$.

\noindent\textbf{Solu\c{c}\~ao:} Dados $(x,y)$, $(z,t) \in \z\times\z$ temos
\begin{align*}
	f((x,y) + (z,t)) &= f(x+z,y+t) = (\overline{x+z - (y+t)},\overline{0}) \\ &= (\overline{x+z - y - t}, \overline{0}) = (\overline{(x-y) + (z-t)},\overline{0}) \\ &= (\overline{x-y},\overline{0}) + (\overline{z-t},\overline{0}) \\ &= f(x,y) + f(z,t).
\end{align*}
Logo $f$ é um homomorfismo de grupos.

Agora
\[
	\ker(f) = \{(x,y) \in \z\times\z \mid f(x,y) = (\overline{0},\overline{0}\}.
\]
Mas
\begin{align*}
	f(x,y) &= (\overline{0},\overline{0})\\
	(\overline{x-y},\overline{0}) &= (\overline{0},\overline{0})\\
	\overline{x-y} &= \overline{0}\\
	\overline{x} &= \overline{y}\\
	x &\equiv y \pmod 7\\
	x &= y + 7k
\end{align*}
para $k \in \z$. Logo
\[
	\ker(f) = \{(y+7k,y) \mid y\in \z,\ k \in \z\}.
\]
\end{document}